\documentclass[11pt,a4paper]{report}
\usepackage[utf-8]{inputenc}
\usepackage[romanian]{babel}
\usepackage{xcolor}
\usepackage{graphicx}
\usepackage{hyperref}
\usepackage{listings}
\usepackage{geometry}
\usepackage{fancyhdr}
\usepackage{setspace}
\usepackage{tocloft}
\usepackage{amsmath}
\usepackage{amssymb}
\usepackage{tikz}
\usepackage{mdframed}

\geometry{left=1.5cm,right=1.5cm,top=2cm,bottom=2cm}
\pagestyle{fancy}
\fancyhf{}
\rhead{Clear2Go — Documentație}
\lhead{\leftmark}
\cfoot{\thepage}

\lstset{
    basicstyle=\ttfamily\small,
    keywordstyle=\color{blue},
    commentstyle=\color{gray},
    stringstyle=\color{red},
    breaklines=true,
    postbreak=\mbox{\textcolor{red}{$\hookrightarrow$}\space},
    language=Java,
    frame=single,
    rulecolor=\color{black!30}
}

\hypersetup{
    colorlinks=true,
    linkcolor=blue!70!black,
    urlcolor=blue!70!black
}

\title{
    \vspace{2cm}
    \Huge \textbf{Clear2Go}\\[0.5em]
    \Large \textit{Platformă de Comunicații Aeronautice}\\[0.3em]
    \normalsize \textit{Documentație Tehnică}\\[2em]
    \large \textbf{\textcolor{blue!70!black}{Prezentată la Faza Națională}}\\[0.3em]
    \normalsize Concurs Național de Informatică / Proiecte IT — 2024
}
\author{
    \textbf{Neculai-Mirea Andrei-Laurențiu} \& \textbf{Dunel Ștefan-Octavian}\\[0.3em]
    \textit{Clasa a X-a A}\\
    Colegiul Național Mihai Viteazul, Slobozia, Ialomița
}
\date{Iulie 2024}

\begin{document}

\maketitle
\thispagestyle{empty}

\vspace{3cm}

\begin{mdframed}[backgroundcolor=blue!5, linecolor=blue!50, linewidth=1.5pt, innertopmargin=10pt, innerbottommargin=10pt]
\begin{center}
\large\textbf{[*] Faza Națională — Proiect Calificat și Prezentat}\\[0.5em]
\normalsize
Clear2Go a fost dezvoltat de doi elevi de clasa a X-a și a ajuns până la
\textbf{faza națională} a competiției, demonstrând că o soluție creată în liceu
poate adresa o problemă reală din aviația civilă românească.
\end{center}
\end{mdframed}

\newpage
\tableofcontents
\newpage

% ============================================================
\chapter{Rezumat Executiv}
% ============================================================

\textbf{Clear2Go} este o platformă inovatoare de comunicații aeronautice care rezolvă
problemele persistente legate de comunicațiile radio dintre piloți și controlorii de
trafic aerian. Aplicația gestionează în timp real autorizațiile emise de turnurile de
control ale aerodromurilor civile, digitalizând un flux de lucru care în mod tradițional
se bazează exclusiv pe radio.

\vspace{0.5em}

\textbf{Obiectiv principal:} Digitalizarea fluxului de autorizare a aeronavelor
(pornire motor, rulaj, intrare aliniere, decolare, aterizare) cu urmărire GPS în timp
real și un sistem de aprobare a cererilor, accesibil atât de pe mobil, cât și de pe web.

\vspace{0.5em}

\begin{mdframed}[backgroundcolor=green!5, linecolor=green!60!black, linewidth=1pt]
\textbf{De ce contează?} \\
Comunicațiile radio în aviație sunt critice — un mesaj ratat sau o autorizare
ambiguă poate genera incidente. Clear2Go adaugă un canal digital redundant, cu
confirmare vizuală clară, reducând riscul de eroare umană.
\end{mdframed}

% ============================================================
\chapter{Context: Competiția Națională}
% ============================================================

\section{Drumul spre faza națională}

Clear2Go a fost conceput, proiectat și implementat integral de doi elevi de clasa a X-a
de la \textbf{Colegiul Național Mihai Viteazul din Slobozia, județul Ialomița}:

\begin{itemize}
    \item \textbf{Neculai-Mirea Andrei-Laurențiu}
    \item \textbf{Dunel Ștefan-Octavian}
\end{itemize}

Proiectul a parcurs etapele de selecție și a ajuns până la \textbf{faza națională} a
competiției, unde a fost prezentat în fața unui juriu de specialiști. Acest parcurs
confirmă că ideea, implementarea și documentația au atins un nivel de calitate
recunoscut la nivel național.

\section{Ce face Clear2Go special în contextul competiției}

\begin{enumerate}
    \item \textbf{Problemă reală, soluție reală} — Nu este un proiect didactic teoretic.
    Clear2Go adresează o problemă identificată în aviația civilă și a fost testat cu
    piloți reali care au oferit feedback pozitiv.

    \item \textbf{Stack tehnologic profesional} — Firebase, OAuth 2.0, Google Maps API,
    Java Android — aceeași tehnologie folosită în aplicații comerciale la scară largă.

    \item \textbf{Dublu front: mobil + web} — Proiectul livrează simultan o aplicație
    Android și un site web cu hartă în timp real, ambele conectate la același backend.

    \item \textbf{Securitate serioasă} — Autentificare OAuth cu Google, criptare
    SHA-256 a transmisiunilor și reguli de securitate Firebase, la standardele industriei.

    \item \textbf{Scalabilitate demonstrată} — Arhitectura suportă teoretic până la
    \textbf{12.500 de aeronave concurente}, cu posibilitate de extindere la mai multe
    aerodromuri.
\end{enumerate}

% ============================================================
\chapter{Arhitectură și Flux de Utilizare}
% ============================================================

\section{Prezentare generală}

\subsection{Roluri în sistem}

\begin{itemize}
    \item \textbf{Pilot} — Solicită autorizații în ordine secvențială pentru fazele
    zborului, cu vizualizare clară a stării fiecărei cereri.
    \item \textbf{Controlor de Trafic Aerian (CTA)} — Aprobă sau respinge cererile
    și monitorizează pozițiile tuturor aeronavelor în timp real pe hartă.
\end{itemize}

\subsection{Fluxul de autorizare}

Cererile pilotului urmează un flux secvențial protejat:
\[
\text{Pornire motor} \rightarrow \text{Rulaj} \rightarrow \text{Intrare aliniere}
\rightarrow \text{Decolare} \rightarrow \text{Aterizare}
\]

Fiecare pas se deblochează doar după aprobarea pasului anterior.
Stările butoanelor sunt vizuale și intuitive:

\begin{itemize}
    \item \textcolor{red}{\textbf{Roșu}} — Stare inițială / Respinsă
    \item \textcolor{orange}{\textbf{Portocaliu}} — În așteptare (cerere trimisă)
    \item \textcolor{green}{\textbf{Verde}} — Aprobată
\end{itemize}

\subsection{Componentele sistemului}

\begin{enumerate}
    \item \textbf{Aplicații mobile} (Android, Java)
    \begin{itemize}
        \item Aplicație pilot (\texttt{FlyActivity.java})
        \item Aplicație controlor (\texttt{ControlActivity.java})
    \end{itemize}

    \item \textbf{Infrastructură backend} (Firebase — Google)
    \begin{itemize}
        \item Realtime Database: starea autorizărilor, coordonate GPS
        \item Firestore: metadate utilizatori, gestionarea rolurilor
        \item Authentication: OAuth Google Sign-In
    \end{itemize}

    \item \textbf{Interfață web} (HTML/JavaScript + Vite)
    \begin{itemize}
        \item Vizualizare în timp real a pozițiilor aeronavelor
        \item Afișarea numelui pilotului la click dreapta pe aeronavă
    \end{itemize}

    \item \textbf{Servicii de localizare}
    \begin{itemize}
        \item GPS nativ al dispozitivului
        \item Google Maps API v2 (vizualizare)
    \end{itemize}
\end{enumerate}

\section{Arhitectura datelor}

\subsection{Structura Firebase Realtime Database}

\begin{lstlisting}
Utilizare/
  Aviatie/
    Aerodromuri/
      AR_AT Bucuresti/
        Flota/
          Avioane/
            [ID_AERONAVA]/
              lat: double
              lng: double
              heading: double
              Pornire motor: boolean
              Rulaj: boolean
              Intrare si aliniere: boolean
              Plecare: boolean
              Aterizare: boolean

Requests/
  [ID_AERONAVA]/
    [TIP_AUTORIZARE]: boolean
\end{lstlisting}

\subsection{Colecții Firestore}

\begin{lstlisting}
Avioane/
  AT Bucuresti/
    Users/
      Admin/[email_admin]
      Pilot/[email_pilot]
      Cz/[email_controlor]
\end{lstlisting}

% ============================================================
\chapter{Securitate}
% ============================================================

\section{Autentificare}

\subsection{OAuth 2.0 cu Google Sign-In}

Clear2Go folosește standardul industrial \textbf{OAuth 2.0} prin Firebase Authentication,
la fel cum procedează aplicațiile enterprise din întreaga lume:

\begin{lstlisting}
GoogleSignInOptions gso = new GoogleSignInOptions.Builder(
    GoogleSignInOptions.DEFAULT_SIGN_IN)
    .requestIdToken(getString(R.string.default_web_client_id))
    .requestEmail()
    .build();
\end{lstlisting}

Fluxul de autentificare:
\begin{enumerate}
    \item Utilizatorul inițiază autentificarea Google
    \item Providerul OAuth Google autentifică
    \item Se obține token-ul ID
    \item Firebase Authentication validează token-ul
    \item UID persistent generat pentru sesiune
    \item Sesiunea este menținută automat de FirebaseAuth
\end{enumerate}

\subsection{Managementul sesiunilor}

\begin{itemize}
    \item Firebase gestionează sesiunea persistent și automat
    \item UID mapat la rolurile utilizatorului în Firestore
    \item La delogare, sesiunea este invalidată complet
\end{itemize}

\section{Criptare și confidențialitate}

\subsection{Securitate în tranzit}

\begin{itemize}
    \item Toate comunicațiile Firebase: criptate \textbf{HTTPS/TLS} (gestionat de Google)
    \item Metadate sensibile: hashing \textbf{SHA-256} înainte de stocare
    \item Coordonatele GPS sunt criptate înainte de trimiterea la Firebase
\end{itemize}

\subsection{Reguli de securitate Firebase}

Accesul la date este restricționat la nivel de bază de date:

\begin{lstlisting}
{
  "rules": {
    "Utilizare": {
      ".read": "auth != null",
      ".write": "auth != null",
      "Aviatie": {
        "Aerodromuri": {
          "$aerodrome": {
            "Flota": {
              "Avioane": {
                ".indexOn": ["lat", "lng", "heading"]
              }
            }
          }
        }
      }
    }
  }
}
\end{lstlisting}

\section{Control acces bazat pe roluri (RBAC)}

Firestore stochează rolurile utilizatorilor, iar aplicația le aplică la nivel de activitate:

\begin{itemize}
    \item \textbf{Admin} — Acces complet
    \item \textbf{Pilot} — Acces la aplicația pilot, generare cereri
    \item \textbf{Cz (Controlor Zboruri)} — Acces la aplicația controlor, aprobare cereri
\end{itemize}

Fiecare activitate Android verifică rolul înainte de a se afișa:

\begin{itemize}
    \item \texttt{MainActivity}: doar autentificare
    \item \texttt{ProfileActivity}: utilizatori autentificați
    \item \texttt{FlyActivity}: exclusiv piloți autorizați
    \item \texttt{ControlActivity}: exclusiv controlori autorizați
\end{itemize}

% ============================================================
\chapter{Scalabilitate și Performanță}
% ============================================================

\section{Capacitate de scalare}

Firebase Realtime Database suportă până la \textbf{25.000 de conexiuni concurente}.
Cu actualizări GPS la 2 Hz per aeronavă:

\[
\text{Aeronave maxime} \approx 12.500
\]

Aceasta depășește cu mult nevoile oricărui aerodrom civil din România.

\section{Performanță aplicație mobilă}

\begin{itemize}
    \item \textbf{Actualizări locație}: 2-5 Hz, consum CPU neglijabil (0–2\%)
    \item \textbf{Randare hartă} (ControlActivity): 16ms/cadru, fluid
    \item \textbf{Memorie de bază}: 20–30 MB
    \item \textbf{Capacitate}: 500 aeronave simultan $\approx$ 120 MB (acceptabil pe
    orice telefon modern cu min. 2 GB RAM)
\end{itemize}

\section{Eficiență rețea}

\begin{itemize}
    \item \textbf{Citiri}: actualizări GPS (5 bytes lat/lng) × 2 Hz × 500 aeronave
    = \textbf{5 KB/sec}
    \item \textbf{Scrieri}: poziție aeronavă (15 bytes) × 2 Hz = 30 bytes/sec
    per aeronavă
    \item Lățime de bandă totală: $\approx$ 200 KB/sec la 500 aeronave — neglijabilă
    pe orice conexiune modernă
\end{itemize}

\subsection{Optimizare scrieri batch}

Aplicația folosește scrieri grupate pentru eficiență maximă:

\begin{lstlisting}
// Scriere batch optimizata
Map<String, Object> updates = new HashMap<>();
updates.put("lat", latitude);
updates.put("lng", longitude);
updates.put("heading", heading);
avionData.updateChildren(updates); // O singura operatiune
\end{lstlisting}

\section{Proiecții scalabilitate}

\begin{tabular}{|l|c|c|}
\hline
\textbf{Metrică} & \textbf{Actual} & \textbf{Limită arhitectură} \\
\hline
Aeronave concurente & 100 & 12.500 \\
Latență query DB & 200–500 ms & $<$ 1s \\
Memorie/aeronavă & 2.5 MB & $<$ 100 MB total \\
\hline
\end{tabular}

\section{Extindere multi-aerodrom}

Arhitectura permite adăugarea oricărui aerodrom prin parametrizarea rutei
în baza de date:

\begin{lstlisting}
// Propunere: selectie aerodrom din ProfileActivity
Intent intent = new Intent(activity, FlyActivity.class);
intent.putExtra("aerodrom", aerodromSelectat);
intent.putExtra("aeronava", aeronavaSelectata);
\end{lstlisting}

% ============================================================
\chapter{Specificații Tehnice}
% ============================================================

\section{Cerințe sistem}

\subsection{Aplicație Android}

\begin{itemize}
    \item \textbf{API minim}: 24 (Android 7.0)
    \item \textbf{API țintă}: 34 (Android 14)
    \item \textbf{Acoperire dispozitive}: 75\% din dispozitivele active
    \item \textbf{RAM minim}: 2 GB
    \item \textbf{Stocare}: 50 MB
    \item \textbf{Rețea}: 3G+ (WiFi recomandat)
\end{itemize}

\subsection{Infrastructură backend}

\begin{itemize}
    \item Firebase Realtime Database (Google-hosted)
    \item Firebase Firestore (Google-hosted)
    \item Firebase Authentication
    \item Google Cloud Platform (GCP)
\end{itemize}

\subsection{Stack de dezvoltare}

\begin{tabular}{|l|l|}
\hline
\textbf{Componentă} & \textbf{Tehnologie} \\
\hline
Limbaj & Java 8 \\
IDE & Android Studio \\
Build System & Gradle 8.x \\
VCS & Git/GitHub \\
Backend & Firebase (Google) \\
Hărți & Google Maps API v2 \\
Frontend web & HTML, JavaScript, Vite \\
\hline
\end{tabular}

\section{Dependințe principale}

\begin{lstlisting}
// Core Android
androidx.appcompat:appcompat:1.7.0
com.google.android.material:material:1.12.0

// Firebase (BOM 33.1.2)
com.google.firebase:firebase-database:21.0.0
com.google.firebase:firebase-auth:23.0.0
com.google.firebase:firebase-firestore:25.0.0
com.google.firebase:firebase-analytics

// Hărti si localizare
com.google.android.gms:play-services-maps:19.0.0
com.google.android.gms:play-services-location:21.3.0
\end{lstlisting}

% ============================================================
\chapter{Foaie de Parcurs}
% ============================================================

\section{Faza 2: Suport multi-aerodrom}

\begin{itemize}
    \item Parametrizarea rutelor de baze de date
    \item Acces bazat pe rol per aerodrom
    \item Dashboard web pentru gestionare multi-site
\end{itemize}

\section{Faza 3: Funcționalități avansate}

\begin{itemize}
    \item Secvențierea predictivă a aeronavelor (bazată pe ML)
    \item Detectarea conflictelor de trafic în timp real
    \item Integrare voce (VoIP backup)
    \item Analize istorice (ore de vârf, blocaje)
\end{itemize}

\section{Faza 4: Deployment enterprise}

\begin{itemize}
    \item Backend on-premises (alternativă la Firebase)
    \item Sisteme de redundanță și failover
    \item Garanții SLA (99.99\% uptime)
    \item Integrare cu sistemele ATC existente
\end{itemize}

% ============================================================
\chapter{Concluzie}
% ============================================================

\section{Rezumat}

Clear2Go este o platformă funcțională de comunicații aeronautice care adresează
\textbf{lacune operaționale reale}. Arhitectura bazată pe Firebase permite
un deployment rapid, scalabil și sigur, adecvat de la MVP până la implementări
de scală medie (500–1.000 de aeronave concurente).

Proiectul a demonstrat că doi elevi de liceu pot construi o soluție tehnologică
completă, cu stack profesional, securitate la standarde industriale și o interfață
utilizator validată de piloți reali — și pot duce această soluție până la
\textbf{faza națională} a unei competiții.

\section{Puncte forte}

\begin{itemize}
    \item \textbf{Autentificare OAuth} (standard industrial)
    \item \textbf{Urmărire GPS în timp real} a tuturor aeronavelor
    \item \textbf{Interfețe intuitive} pilot / controlor, validate cu utilizatori reali
    \item \textbf{Arhitectură modulară} ușor extensibilă
    \item \textbf{Dublu frontend}: mobil (Android) + web (Vite/JS)
    \item \textbf{Versionare Git/GitHub} cu istoricul complet al dezvoltării
    \item \textbf{Scalabilitate teoretică}: 12.500+ aeronave concurente
\end{itemize}

\section{Impactul potențial}

\[
\text{MVP} \rightarrow \text{Alpha — Aerodrom unic} \rightarrow
\text{Beta — Multi-aerodrom} \rightarrow \text{Producție Națională}
\]

\vspace{1em}

\begin{mdframed}[backgroundcolor=blue!5, linecolor=blue!50, linewidth=1.5pt]
\begin{center}
\textbf{Clear2Go — dezvoltat de elevi, validat național, gândit pentru viitor.}\\[0.3em]
\small Neculai-Mirea Andrei-Laurențiu \& Dunel Ștefan-Octavian\\
Colegiul Național Mihai Viteazul, Slobozia · Clasa a X-a A · 2024
\end{center}
\end{mdframed}

\vspace{2cm}
\noindent \textit{Versiunea documentului: 1.0 — Faza Națională}\\
\textit{Data: Iulie 2024}\\
\textit{Clasificare: Documentație Tehnică (Non-Confidențial)}

\end{document}
